% English draft based on appendices/appendix_d_optional_notes_ja.tex
\section{Optional Notes}

\noindent\textbf{Scope of the account.} These optional notes and, in the Japanese manuscript, argument sketches retain their different documentary roles. The sketches are not certified proofs. Strict-increase premises, stopping-as-fixed-point interpretation, barrier inaccessibility, emergence, field/network identity, and mobility/error equivalence remain separate burdens. The AI-cost comparison is a possible reading, not a derived cost necessity or a tested computational result.


This appendix records optional interpretations that are useful but not necessary for the main argument.

\subsection{Psychological Projection}

The structural projection used in this paper can resemble psychological projection, because internally unavailable stopping is projected into external or shared loci. However, the present theory does not presuppose subject, unconscious mechanism, or clinical defense. It remains a structural account.

\subsection{Science}

Science is not treated as a simple victory over religion. It is the institutionalization of revisable stopping. It maintains stopping while also building procedures for reopening, testing, and replacing stops.

\subsection{AI Cost}

The qualitative phase proposal motivates comparing inference cost with possible logging, external-constraint, and stopping costs. It does not derive a necessary bandwidth or cost law, and no measured cost relation is supplied.

\subsection{Freedom}

Freedom is not the absence of stopping. Within the proposed functional interpretation, movable stopping may support free inference; a universal or phenomenal necessity is not established. Dissolution without stopping and hardening without mobility are qualitative interpretations within that proposal, not universal laws.

\bigskip
